%----------------------------------------------------------------------------------------------------
% START OF CHAPTER
%----------------------------------------------------------------------------------------------------
\chapter{Apr 11}
\paragraph{Topics} 
\textit{Numbers. Time. \iti{-are} verb form recap. Sentence review.}

\section{Grammar} \label{0411-Grammar}
\subsection{Numbers}
See Table \ref{tab:tNumbersZeroToTwenty}. Numbers are one of those things you just have to learn. 

\begin{table}[ht] %[ht] % place table roughly [h]ere or else at the [t]op of page
	\centering
	\begin{tabular}{ r l | r l | r l }
		\textbf{\#} &              & \textbf{\#} &                  & \textbf{\#} &                     \\
		\hline              
		\textbf{0} & \iti{zero}    & \textbf{7}  & \textit{sette}   & \textbf{14} & \textit{quattordici} \\
		\textbf{1} & \iti{uno}     & \textbf{8}  & \textit{otto}    & \textbf{15} & \textit{quindici}    \\
		\textbf{2} & \iti{due}     & \textbf{9}  & \textit{nove}    & \textbf{16} & \textit{sedici}      \\
		\textbf{3} & \iti{tre}     & \textbf{10} & \textit{dieci}   & \textbf{17} & \textit{diciassette} \\
		\textbf{4} & \iti{quattro} & \textbf{11} & \textit{undici}  & \textbf{18} & \textit{diciotto}    \\
		\textbf{5} & \iti{cinque}  & \textbf{12} & \textit{dodici}  & \textbf{19} & \textit{diciannove} \\
		\textbf{6} & \iti{sei}     & \textbf{13} & \textit{tredici} & \textbf{20} & \textit{venti} 
	\end{tabular}
	\caption{Numbers zero through twenty}
	\label{tab:tNumbersZeroToTwenty}
\end{table}

Once the base forms are learned, e.g., \iti{venti}, `twenty', compound forms like 
`twenty-one' are easy to figure out. 

\begin{table}[ht] %[ht] % place table roughly [h]ere or else at the [t]op of page
	\centering
	\begin{tabular}{ r l | r l  }
		\textbf{\#} &                 & \textbf{\#}   &                \\
		\hline              
		\textbf{10} & \iti{dieci}     & \textbf{70}   & \iti{settanta} \\
		\textbf{20} & \iti{venti}     & \textbf{80}   & \iti{ottanta}  \\
		\textbf{30} & \iti{trenta}    & \textbf{90}   & \iti{novanta}  \\
		\textbf{40} & \iti{quaranta}  & \textbf{100}  & \iti{cento}    \\
		\textbf{50} & \iti{cinquanta} & \textbf{1000} & \iti{mille}    \\
		\textbf{60} & \iti{sessanta}  &               &                
	\end{tabular}
	\caption{Multiples of ten and one hundred}
	\label{tab:tNumbersTenToOneThousand}
\end{table}

For large numbers (thousands), you would speak the thousands part, then \iti{e}, and then the remainder.
Phone numbers can be read out number by number: `seven-zero-nine' / \iti{`sette-zero-nove'} style.

\paragraph{Examples}
\begin{itemize}
  \item{Twenty-five} \\
    \iti{Venticinque}
  \item{One hundred and twenty-five} \\
    \iti{Centoventicinque}
  \item{\euro{2350}}  \\
    \iti{Due milla e trecentocinquanta euro.}
\end{itemize}

Note: The word `euro' is pronounced like \textit{Eh-OO-roh}, not \textit{YOO-roh}.

\subsection{Time}
\textmd{Clocks in Italy use the 24-hour system for official business, such as train times, 
  although speakers might use expressions like \iti{di sera} in casual conversation. 
  So for example a train announcer might say \iti{Il treno per Bologna parte alle quindici} 
  but you might refer to the time as \iti{tre di sera}.
	Notice how \iti{di sera} spans everything after noon, and not just the evening time.}
	
\paragraph{Examples}
\begin{itemize}
  \item{Six \textsc{PM}} \\
    \iti{Diciotto.} or \iti{Sei di sera.}
  \item{Six \textsc{AM}} \\
    \iti{Sei di mattina.}
  \item{5:45 \textsc{AM}} \\
    \iti{Sei meno un quarto.}
  \item{7:05 \textsc{AM}} \\
    \iti{Sette e cinque.}
  \item{8:30 \textsc{AM}} \\
    \iti{Otto e mezzo.}
  \item{Noon} \\
    \iti{Il mezzogiorno.}
  \item{13:55} \\
    \iti{Tredici e cinquantacinque.}
\end{itemize}

\subsection{Verb recap: the regular \iti{-are} form}
See Table \ref{tab:tPortare} as an example. The regular verbs ending \iti{-are} can be conjugated,
in their present tense, by stripping off the \iti{-are} ending and adding \iti{-o, -i, -a, -iamo, -ate, -ano}.
For example, \iti{io porto}, `I bring' from \iti{portare}, `to bring, to deliver'.

\begin{table}[ht] %[ht] % place table roughly [h]ere or else at the [t]op of page
	\centering
	\begin{tabular}{ r | l }
		\itb{portare} & \textit{to bring} \\
		\hline              
		\itb{io}      & \iti{porto}    \\
		\itb{tu}      & \iti{porti}    \\
		\itb{lui/lei} & \iti{porta}    \\
		\itb{noi}     & \iti{portiamo} \\
		\itb{voi}     & \iti{portate}  \\
		\itb{loro}    & \iti{portano}  
	\end{tabular}
	\caption{An example of a regular \iti{-are} verb: \iti{portare}}
	\label{tab:tPortare}
\end{table}

\subsection{Verb recap: the regular \iti{-ire} form}

See Table \ref{tab:tAprire} as an example. The regular verbs ending \iti{-ire} can be conjugated,
in their present tense, by stripping off the \iti{-ire} ending and adding \iti{-o, -i, -e, -iamo, -ite, -ono}.
For example, \iti{io apro}, `I open' from \iti{aprire}, `to open'. A second group of \iti{-ire} verbs,
not covered here, uses endings like \iti{-isco}, e.g., \iti{io capisco} from \iti{capire}, `to understand'.
See \cite{Lawless}.

\begin{table}[ht] %[ht] % place table roughly [h]ere or else at the [t]op of page
	\centering
	\begin{tabular}{ r | l }
		\itb{aprire} & \textit{to open} \\
		\hline              
		\itb{io}      & \iti{apro}    \\
		\itb{tu}      & \iti{apri}    \\
		\itb{lui/lei} & \iti{apre}    \\
		\itb{noi}     & \iti{apriamo} \\
		\itb{voi}     & \iti{aprite}  \\
		\itb{loro}    & \iti{aprono}  
	\end{tabular}
	\caption{An example of a regular \iti{-ire} verb: \iti{aprire}}
	\label{tab:tAprire}
\end{table}

\section{Vocabulary} \label{0411-Vocabulary}
\subsection{Words}

Very varied list of words this week - Table \ref{tab:tVocabulary-0411}.

\begin{table}[ht] %[ht] % place table roughly [h]ere or else at the [t]op of page
	\centering
	\begin{tabular}{ r l  | r l  }
		to bring       & \iti{portare}     & to go                & \iti{andare} \\
		dinner         & \iti{cena}        & cream                & \iti{crema}  \\
		pistachio      & \iti{pistacchio}  &                      &              \\
		               &                   &                      &              \\
		you (formal)   & \iti{Lei}         & excuse me (formal)   & \iti{scusi}  \\
		you (informal) & \iti{tu}          & excuse me (informal) & \iti{scusa}  \\
	                   &                   &                      &              \\
        How are you? (formal)   & \iti{Come sta?}  &              &              \\
        How are you? (informal) & \iti{Come stai?} &              &              \\
	                   &                   &                      &              \\
        First time     & \iti{Prima volta} & Friends              & \iti{amici}  \\
        Cash           & \iti{contanti}    & Credit card          & \iti{carta di credito} \\
	                   &                   &                      &              \\
	    Closed         & \iti{chiuso}      & Open                 & \iti{aperto} \\
	    To open        & \iti{aprire}      & To Push/pull         & \iti{spingere/tirare}
	\end{tabular}
	\caption{Vocab list for Apr 11}
	\label{tab:tVocabulary-0411}
\end{table}

Watch out for this `false friend' word pair:

\begin{enumerate}
  \item `Parents' is \iti{genitori}
  \item `Relatives' is \iti{parenti}
\end{enumerate}

\paragraph{Examples}
\begin{itemize}
	\item I go, you go, we go, you (pl.) go.     \\
	\iti{Io vado, tu vai, noi andiamo, voi andate.}
	
	\item Do you have wine? (formal)      \\
    \iti{Lei ha vino?}
		
	\item Do you have wine? (informal)    \\
	\iti{Voi avete vino?}
	
    \item I do not understand.            \\
    \iti{Non ho capisco.}

    \item Talk slow(ly).                  \\
    \iti{Parla piano.}

    \item Who's paying?                   \\
    \iti{Chi paga?}	
\end{itemize}

\section{Conversation - Sentence review}
\begin{itemize}
  \item I'm talking with Claudio                   \\
    \iti{Io sto parlando con Claudio.}
  \item I'm at home now.                           \\
    \iti{Io sono a casa adesso.}
  \item I'm staying home.                          \\
    \iti{Io sto a casa.}
  \item I would like a bottle of red wine, please. \\
    \iti{Io vorrei una bottiglia di vino rosso, per favore.}
  \item I would eat \ldots                         \\
    \iti{Io mangerei \ldots}
  \item I am eating.                               \\
    \iti{Io sto mangiando.}
  \item I'm tired. (Male speaker)                  \\
    \iti{Io sono stanco.}
  \item I am.                                      \\
    \iti{Io sono.}
  \item I am speaking with Claudio.                \\
    \iti{Io sto parlando con Claudio.}
  \item I would like to pay the bill.              \\
    \iti{Io vorrei pagare il conto.}
  \item My name is L, pleased to meet you, I am American. \\
    \iti{Il mio nome è L, piacere, sono americana.}
  \item I have cash.                               \\
    \iti{Io ho contanti.}
  \item I don't want this, thank you.              \\
    \iti{Io non voglio questo, grazie.}
  \item I can pay.                                 \\
    \iti{Io posso pagare.}
  \item I'm driving to Rome.                       \\
    \iti{Io sto guidando a Roma.}
  \item I'm not tired today. (Female speaker)      \\
    \iti{Io non sono stanca oggi.}
  \item I'm going to the store this evening.       \\
    \iti{Io sto andando al negozio questa sera.}
  \item I'm going home because I'm very tired. (Male speaker) \\
    \iti{Io vado a casa perché io sono molto stanco.}
  \item Let's go!                                  \\
    \iti{Andiamo!}
  \item Letś eat!                                  \\
    \iti{Mangiamo!}
  \item We're happy because we are studying Italian. \\
    \iti{Noi siamo felici perché noi stiamo studiando italiano}\footnote{Notice \iti{felici}, plural of \iti{felice}, to match the plurality of the speakers (`we').}
  \item I would (like to) try that dish, please.   \\
    \iti{Io proverei questo piatto, per favore.}
\end{itemize}

\subsection{Gestures - video}
See the video at \cite{Virgallito}
%----------------------------------------------------------------------------------------------------
% END OF CHAPTER
%----------------------------------------------------------------------------------------------------
